\chapter{ZAKLJUČAK}
\label{ch:zakljucak}

U radu je razvijen i eksperimentalno vrednovan sustav za autonomnu interakciju mobilnog
manipulatora s okolinom, na primjeru otvaranja dvaju vrsti vrata. Sustav objedinjuje percepciju,
procjenu sile i momenta te upravljanje u kontaktu, a zadatak upravljanja riješen je dvama
pristupima: klasičnim, temeljenim na modelu, i pristupom temeljenim na podržanom učenju.
Pristupi su potom uspoređeni na istom robotu, istoj sceni i istim mjerilima.

Modul percepcije određuje položaj i orijentaciju kvake iz slike, pomoću vizualnih oznaka.
Postupak procjene sile i momenta na vrhu alata izveden je iz momenata u zglobovima, bez zasebnog
senzora sile, a iz odnosa procijenjene sile i ostvarenog pomaka pokazano je kako se procjenjuju
krutost okoline i geometrija ograničenja gibanja.

Klasični pristup rješava cijeli zadatak, od prilaska i hvata preko otvaranja do prolaska kroz
vrata, pri čemu tip vrata prepoznaje sam, iz geometrije kvake. Pristup temeljen na učenju
formuliran je kao Markovljev proces odlučivanja s varijabilnom impedancijom u prostoru akcije,
čime politika u svakom koraku bira koliko će regulator biti krut po svakoj osi.

\section{Glavni nalazi}

Usporedba dvaju pristupa dala je rezultat koji se ne svodi na prednost jednoga nad drugim, nego
ovisi o tipu vrata, i upravo je ta ovisnost najvrjedniji nalaz rada.

Kod kliznih vrata oba pristupa dosežu puni cilj u svakom pokušaju. Ograničenje je ondje pravac,
smjer sile poklapa se sa smjerom gibanja, i eksplicitno poznavanje geometrije daje sve što je
potrebno, pa dodatna sloboda u odabiru krutosti ne donosi mjerljivu prednost.

Kod zakretnih vrata razlika je velika. Klasični pristup pouzdano otvara vrata do 50\textdegree{},
ali preko otprilike 57\textdegree{} ne može, jer kvaka klizne iz prihvatnice. Naučena politika
vrata otvara preko 90\textdegree{} u 96 \% epizoda, s medijanom otvora na 122\textdegree{} i
gornjim kvartilom praktički na mehaničkoj granici zgloba. Uz to, klasični pristup pri otvaranju
zakretnih vrata redovito prelazi 400 N, s vršnom silom od 625 N, dok naučena politika tu granicu
prekoračuje u manje od 1 \% epizoda.

Zaključak je time uvjetovan, a ne općenit: eksplicitno poznavanje geometrije dovoljno je kad je
geometrija jednostavna i dobro procijenjena, a gubi prednost čim zadatak zahtijeva prilagodbu
krutosti tijekom same interakcije.

Dva pristupa razlikuju se i po tome što o zadatku znaju. Klasični pristup položaj osi šarke
izračunava iz poze oznake i poznate širine krila, uz pogrešku od 5 do 8 mm, dok naučena politika
tu geometriju nema ni u kojem obliku i mora je otkriti iz vlastitih opažanja sile i pomaka. S
druge strane, naučena politika koordinaciju platforme i manipulatora uči, dok je klasični pristup
rješava unaprijed propisanim slijedom faza. Upravo ta propisana shema kod zakretnih vrata
postaje ograničenje.

\section{Ograničenja}

Cjelokupan je rad proveden u simulaciji, pa se zaključci odnose na model, a ne izravno na stvarni
robot. Uz to, sustav sadrži nekoliko svjesnih pojednostavljenja i nekoliko neriješenih pitanja.

Zglob koji bi omogućio zakretanje kvake zavaren je u oba pristupa, čime je zadatak sveden na
jedan stupanj slobode po tipu vrata. Percepcija se oslanja na vizualne oznake postavljene na
vrata, pa sustav u nepripremljenoj okolini ne bi radio. Mecanum kotači nisu simulirani kao tijela
u kontaktu s podlogom, nego je gibanje platforme zadano kinematički.

Platforma ostvaruje tek dio naređene brzine, a uzrok tome nije utvrđen unatoč tome što je više
mogućih objašnjenja provjereno i isključeno. Posljedica je da se prijeđeni put platforme ne smije
uzimati kao pouzdana mjera gibanja, zbog čega provjera napretka pri otvaranju kliznih vrata čita
stanje zgloba izravno iz simulatora. Ovo ograničenje vrlo je vjerojatno posljedica simulatora,
i vjerojatno ne bi bila problem na stvarnom robotu.

Naposljetku, sustav sadrži postupak procjene geometrije ograničenja iz odnosa sile i gibanja, ali
on nije povezan s izvođenjem: u klasičnom se pristupu os šarke izračunava iz percepcije, a
procjena iz gibanja postoji kao zaseban, samostalno provjeren postupak.

\section{Smjernice za daljnji rad}

\begin{itemize}
    \item \textbf{Prepozicioniranje hvata.} Jedini put preko granice od 57\textdegree{} kod
          zakretnih vrata koji ne traži bočni prostor kojega u sceni nema jest da manipulator
          pusti kvaku, platforma se premjesti, a kvaka se ponovno uhvati iz povoljnije
          konfiguracije.

    \item \textbf{Povezivanje procjene ograničenja s izvođenjem.} Pokretanje procjene odmah
          nakon hvata i odabir načina otvaranja na temelju njezina ishoda uklonilo bi oslanjanje
          na tip kvake prepoznat iz percepcije i učinilo sustav neovisnim o poznatoj geometriji.

    \item \textbf{Odometrija platforme iz percepcije.} Riješila bi oslanjanje na stanje zgloba
          očitano iz simulatora i omogućila mjerenje prijeđenog puta veličinama dostupnima i na
          stvarnom robotu.

    \item \textbf{Prolazak kroz zakretna vrata.} Kod kliznih vrata platforma ostaje uz istu
          stranu zida, pa je otvor dostupan bočnim gibanjem; kod zakretnih vrata robot samim otvaranjem
          prolazi kroz vrata.

    \item \textbf{Zakretanje kvake.} Vraćanje tog stupnja slobode proširilo bi zadatak na vrata
          sa stvarnom bravom, ali bi tražilo upravljanje koje kut zakreta kvake ne prati kruto.

    \item \textbf{Detekcija bez oznaka.} Zamjena vizualnih oznaka postupkom koji kvaku
          prepoznaje izravno iz slike uklonila bi najizraženije pojednostavljenje u radu, pri
          čemu ostatak sustava ne bi trebalo mijenjati, jer je sučelje prema izvođenju zadatka
          svedeno na pozu hvatišta.

    \item \textbf{Prijenos na stvarni robot.} Konačna provjera svih zaključaka traži izvođenje
          izvan simulacije, pri čemu su procjena sile bez zasebnog senzora i pojednostavljeni
          model gibanja platforme najizgledniji izvori odstupanja.
\end{itemize}